\chapter{Catching Up with Friends at a Cozy Autumn Café}

% ================= PERCAKAPAN 1-5 =================
\begin{convobox}{1}
    \noindent
    \jpword{ひさ}{久}{Hisa}\jpkana{しぶり！}{shiburi!} 
    \jpword{げんき}{元気}{Genki}\jpkana{だった？}{datta?}

    \vspace{1.5em}
    \noindent{\Large \color{articolor} \textbf{Arti:} Lama tak jumpa! Apa kabar?}
\end{convobox}

\begin{convobox}{2}
    \noindent
    \jpkana{うん、}{Un,} 
    \jpword{げんき}{元気}{genki}\jpkana{だったよ。}{datta yo.} 
    \jpkana{たなか}{Tanaka}\jpkana{さんは？}{san wa?}

    \vspace{1.5em}
    \noindent{\Large \color{articolor} \textbf{Arti:} Ya, aku baik-baik saja. Kalau Tanaka-san?}
\end{convobox}

\begin{convobox}{3}
    \noindent
    \jpword{さいご}{最後}{Saigo} \jpkana{に}{ni} 
    \jpword{あ}{会}{a}\jpkana{ったのは、}{tta no wa,} 
    \jpkana{いつだったっけ？}{itsu datta kke?}

    \vspace{1.5em}
    \noindent{\Large \color{articolor} \textbf{Arti:} Terakhir kali kita bertemu, kapan ya?}
\end{convobox}

\begin{convobox}{4}
    \noindent
    \jpkana{えーっと…}{Eetto...} 
    \jpword{きょねん}{去年}{kyonen} \jpkana{の}{no} 
    \jpword{ふゆ}{冬}{fuyu} \jpkana{じゃない？}{ja nai?}

    \vspace{1.5em}
    \noindent{\Large \color{articolor} \textbf{Arti:} Emm... bukannya musim dingin tahun lalu?}
\end{convobox}

\begin{convobox}{5}
    \noindent
    \jpkana{いっしょに}{Issho ni} 
    \jpword{おんせん}{温泉}{onsen} \jpkana{に}{ni} 
    \jpword{い}{行}{i}\jpkana{ったよね。}{tta yo ne.}

    \vspace{1.5em}
    \noindent{\Large \color{articolor} \textbf{Arti:} Kita pergi ke pemandian air panas (onsen) bareng, kan.}
\end{convobox}

% --- LESSON BUFFER 1 ---
\begin{lessonbox}{Catatan Belajar: Selamat Datang di Bahasa Kasual!}
    Teman-teman sadar nggak? Di bab ini nggak ada kata \textbf{"desu"} atau \textbf{"masu"}! Kalau kita ngobrol dengan teman dekat, kita pakai bentuk kasual (Biasa).\\
    
    Misalnya, nanya kabar nggak lagi pakai \textit{"Genki desu ka?"}, tapi cukup \textbf{"Genki datta?"}. Trus, "Ya" bukan lagi \textit{"Hai"}, melainkan \textbf{"Un"}. Lebih santai dan akrab banget, kan?
\end{lessonbox}

% ================= PERCAKAPAN 6-11 =================
\begin{convobox}{6}
    \noindent
    \jpword{さいきん}{最近}{Saikin} \jpkana{は}{wa} 
    \jpkana{どう？}{dou?} 
    \jpword{いそが}{忙}{isoga}\jpkana{しい？}{shii?}

    \vspace{1.5em}
    \noindent{\Large \color{articolor} \textbf{Arti:} Akhir-akhir ini gimana? Sibuk?}
\end{convobox}

\begin{convobox}{7}
    \noindent
    \jpkana{うん、}{Un,} 
    \jpword{しごと}{仕事}{shigoto} \jpkana{が}{ga} 
    \jpkana{けっこう}{kekkou} 
    \jpword{いそが}{忙}{isoga}\jpkana{しい。}{shii.}

    \vspace{1.5em}
    \noindent{\Large \color{articolor} \textbf{Arti:} Ya, kerjaan lumayan sibuk.}
\end{convobox}

\begin{convobox}{8}
    \noindent
    \jpkana{でも}{Demo} 
    \jpword{たの}{楽}{tano}\jpkana{しいよ。}{shii yo.}

    \vspace{1.5em}
    \noindent{\Large \color{articolor} \textbf{Arti:} Tapi menyenangkan kok.}
\end{convobox}

\begin{convobox}{9}
    \noindent
    \jpkana{そっか。}{Sokka.} 
    \jpword{やす}{休}{yasu}\jpkana{みの}{mi no} 
    \jpword{ひ}{日}{hi} \jpkana{は}{wa} 
    \jpkana{なにしてる？}{nani shiteru?}

    \vspace{1.5em}
    \noindent{\Large \color{articolor} \textbf{Arti:} Begitu ya. Pas hari libur ngapain aja?}
\end{convobox}

\begin{convobox}{10}
    \noindent
    \jpword{りょうり}{料理}{Ryouri} \jpkana{を}{o} 
    \jpkana{したり、}{shitari,} 
    \jpkana{カフェ}{kafe} \jpkana{に}{ni} 
    \jpword{い}{行}{i}\jpkana{ったり、}{ttari,}

    \vspace{1.5em}
    \noindent{\Large \color{articolor} \textbf{Arti:} Memasak, pergi ke kafe,}
\end{convobox}

\begin{convobox}{11}
    \noindent
    \jpword{いえ}{家}{Ie} \jpkana{で}{de} 
    \jpkana{ごろごろしたり}{gorogoro shitari} 
    \jpkana{してるよ。}{shiteru yo.}

    \vspace{1.5em}
    \noindent{\Large \color{articolor} \textbf{Arti:} bermalas-malasan di rumah, gitu deh.}
\end{convobox}

% --- LESSON BUFFER 2 ---
\begin{lessonbox}{Catatan Belajar: Ngapain Aja? Pakai \textit{~tari}!}
    Kalau kamu mau cerita beberapa aktivitas yang kamu lakuin (nggak harus berurutan), kamu bisa pakai pola \textbf{~tari, ~tari shimasu}. \\
    
    Contoh di atas: \textit{Ryouri o shi\textbf{tari}} (masak), \textit{Kafe ni i\textbf{ttari}} (pergi ke kafe). Ini asik banget dipakai kalau ditanya "Liburan ngapain aja?"
\end{lessonbox}

% ================= PERCAKAPAN 12-19 =================
\begin{convobox}{12}
    \noindent
    \jpword{さいきん}{最近}{Saikin}\jpkana{、}{,} 
    \jpkana{けっこう}{kekkou} 
    \jpword{さむ}{寒}{samu}\jpkana{くなってきたよね。}{ku natte kita yo ne.}

    \vspace{1.5em}
    \noindent{\Large \color{articolor} \textbf{Arti:} Akhir-akhir ini, (cuacanya) sudah lumayan dingin ya.}
\end{convobox}

\begin{convobox}{13}
    \noindent
    \jpword{さむ}{寒}{Samu}\jpkana{いのは}{i no wa} 
    \jpword{にがて}{苦手}{nigate}\jpkana{？}{?}

    \vspace{1.5em}
    \noindent{\Large \color{articolor} \textbf{Arti:} Kamu nggak tahan sama dingin?}
\end{convobox}

\begin{convobox}{14}
    \noindent
    \jpword{さむ}{寒}{Samu}\jpkana{いのは}{i no wa} 
    \jpword{だいじょうぶ}{大丈夫}{daijoubu}\jpkana{。}{.}

    \vspace{1.5em}
    \noindent{\Large \color{articolor} \textbf{Arti:} Kalau dingin nggak apa-apa (tahan).}
\end{convobox}

\begin{convobox}{15}
    \noindent
    \jpword{あつ}{暑}{Atsu}\jpkana{い}{i} 
    \jpword{ほう}{方}{hou} \jpkana{が}{ga} 
    \jpword{にがて}{苦手}{nigate} \jpkana{かな。}{kana.}

    \vspace{1.5em}
    \noindent{\Large \color{articolor} \textbf{Arti:} Kayaknya aku lebih nggak tahan panas deh.}
\end{convobox}

\begin{convobox}{16}
    \noindent
    \jpword{さいきん}{最近}{Saikin}\jpkana{、}{,} 
    \jpkana{どこか}{doko ka} 
    \jpword{りょこう}{旅行}{ryokou} \jpkana{に}{ni} 
    \jpword{い}{行}{i}\jpkana{った？}{tta?}

    \vspace{1.5em}
    \noindent{\Large \color{articolor} \textbf{Arti:} Akhir-akhir ini, ada pergi liburan ke suatu tempat?}
\end{convobox}

\begin{convobox}{17}
    \noindent
    \jpword{せんげつ}{先月}{Sengetsu}\jpkana{、}{,} 
    \jpword{なごや}{名古屋}{Nagoya} \jpkana{に}{ni} 
    \jpword{い}{行}{i}\jpkana{ったよ。}{tta yo.}

    \vspace{1.5em}
    \noindent{\Large \color{articolor} \textbf{Arti:} Bulan lalu, (aku) pergi ke Nagoya lho.}
\end{convobox}

\begin{convobox}{18}
    \noindent
    \jpkana{ジブリパーク}{Jiburi paaku} \jpkana{と}{to} 
    \jpword{なごやじょう}{名古屋城}{Nagoya-jou} \jpkana{に}{ni} 
    \jpword{い}{行}{i}\jpkana{った。}{tta.}

    \vspace{1.5em}
    \noindent{\Large \color{articolor} \textbf{Arti:} Pergi ke Ghibli Park dan Kastil Nagoya.}
\end{convobox}

\begin{convobox}{19}
    \noindent
    \jpkana{おいしい}{Oishii} 
    \jpword{もの}{物}{mono} \jpkana{も}{mo} 
    \jpkana{たくさん}{takusan} 
    \jpword{た}{食}{ta}\jpkana{べて、}{bete,} 
    \jpkana{すごく}{sugoku} 
    \jpword{たの}{楽}{tano}\jpkana{しかった。}{shikatta.}

    \vspace{1.5em}
    \noindent{\Large \color{articolor} \textbf{Arti:} Makan banyak makanan enak, (pokoknya) sangat menyenangkan.}
\end{convobox}

% --- LESSON BUFFER 3 ---
\begin{lessonbox}{Catatan Belajar: Tahan atau Nggak Tahan?}
    Orang Jepang suka nanya apakah kita bisa tahan sesuatu (makanan pedas, cuaca panas/dingin).
    \begin{itemize}
        \item \textbf{苦手 (Nigate)} = Lemah terhadap sesuatu / Nggak tahan / Kurang suka.
        \item \textbf{大丈夫 (Daijoubu)} = Nggak apa-apa / Aman / Tahan.
    \end{itemize}
    Jadi kalau ditanya "Samui no wa nigate?" artinya "Kamu nggak kuat dingin ya?".
\end{lessonbox}

% ================= PERCAKAPAN 20-25 =================
\begin{convobox}{20}
    \noindent
    \jpkana{あっ、}{A,,} 
    \jpword{きょう}{今日}{kyou} \jpkana{は}{wa} 
    \jpkana{そろそろ}{sorosoro} 
    \jpword{かえ}{帰}{kae}\jpkana{らなくちゃ。}{ranakucha.}

    \vspace{1.5em}
    \noindent{\Large \color{articolor} \textbf{Arti:} Ah, hari ini harus segera pulang nih.}
\end{convobox}

\begin{convobox}{21}
    \noindent
    \jpkana{よかったら、}{Yokattara,} 
    \jpword{らいげつ}{来月}{raigetsu} \jpkana{も}{mo} 
    \jpkana{また}{mata} 
    \jpkana{ランチ}{ranchi} \jpkana{か}{ka} 
    \jpword{おちゃ}{お茶}{o-cha} \jpkana{でも}{demo} 
    \jpkana{しない？}{shinai?}

    \vspace{1.5em}
    \noindent{\Large \color{articolor} \textbf{Arti:} Kalau berkenan, bulan depan mau makan siang atau minum teh lagi nggak?}
\end{convobox}

\begin{convobox}{22}
    \noindent
    \jpkana{いいね！}{Ii ne!}

    \vspace{1.5em}
    \noindent{\Large \color{articolor} \textbf{Arti:} Ide bagus!}
\end{convobox}

\begin{convobox}{23}
    \noindent
    \jpword{へいじつ}{平日}{Heijitsu} \jpkana{は}{wa} 
    \jpword{しごと}{仕事}{shigoto} \jpkana{だけど、}{dakedo,} 
    \jpword{どにち}{土日}{donichi} \jpkana{は}{wa} 
    \jpword{だいたい}{大体}{daitai} 
    \jpword{あ}{空}{a}\jpkana{いてるよ。}{iteru yo.}

    \vspace{1.5em}
    \noindent{\Large \color{articolor} \textbf{Arti:} Kalau hari biasa sih kerja, tapi akhir pekan biasanya kosong (jadwalnya) kok.}
\end{convobox}

\begin{convobox}{24}
    \noindent
    \jpkana{じゃあ、}{Jaa,} \jpkana{また}{mata} 
    \jpword{れんらく}{連絡}{renraku} \jpkana{するね。}{suru ne.}

    \vspace{1.5em}
    \noindent{\Large \color{articolor} \textbf{Arti:} Kalau gitu, nanti aku hubungi lagi ya.}
\end{convobox}

\begin{convobox}{25}
    \noindent
    \jpword{きょう}{今日}{Kyou} \jpkana{は}{wa} \jpkana{ありがとう。}{arigatou.}

    \vspace{1.5em}
    \noindent{\Large \color{articolor} \textbf{Arti:} Terima kasih ya untuk hari ini.}
\end{convobox}

% --- SUMMARY BOX ---
\vspace{2em}
\begin{tcolorbox}[
    colback=blue!5!white,
    colframe=blue!80!black,
    boxrule=3pt,
    arc=5mm,
    title={\huge\bfseries \sffamily 🌟 Rangkuman Bab 3 🌟},
    fonttitle=\color{white},
    attach boxed title to top center={yshift=-4mm},
    boxed title style={colback=blue!80!black, rounded corners},
    left=5mm, right=5mm, top=7mm, bottom=5mm
]
\Large \textbf{Keren banget! Sekarang kamu bisa ngobrol santai!} Di bab ini kita belajar:
\begin{itemize}
    \item \textbf{Sapaan Akrab:} \textit{Hisashiburi!} (Lama tak jumpa!)
    \item \textbf{Menyebut Berbagai Aktivitas:} \textit{\dots shitari, \dots ittari} (Bisa berarti: melakukan X, pergi ke Y, dll).
    \item \textbf{Menyatakan Kesukaan/Kelemahan:} \textit{Nigate} (Kurang suka / Nggak tahan) dan \textit{Daijoubu} (Tahan / Aman).
    \item \textbf{Harus melakukan sesuatu:} \textit{Kaeranakucha!} (Harus pulang!).
    \item \textbf{Berpisah dengan teman:} \textit{Mata renraku suru ne} (Nanti aku kabari / hubungi lagi ya).
\end{itemize}
Coba deh praktekkan ngobrol dengan gaya kasual ini ke teman belajarmu. Pasti asik banget!
\end{tcolorbox}

\newpage